\section{The \texttt{MCViNE\_SQ} McStas Component}
Elastic isotropic scatterer with structure factor S(|Q|) (expression or table).

\subsection*{Identification}
\begin{itemize}
  \item \textbf{Author:} Fahima Islam (McStas port of the MCViNE SQkernel; MCViNE by J. Y. Y. Lin et al.)
  \item \textbf{Origin:} MCViNE, https://github.com/mcvine/mcvine (mccomponents/lib/kernels/sample/SQkernel.cc (+ GridSQ, SQ\_fromexpression))
  \item \textbf{Date:} 2026-09-29
\end{itemize}

\subsection*{Description}
Elastic scattering with angular distribution given by S(|Q|), for liquids, glasses or powder-averaged diffuse scattering. S(Q) is either an expression in Q or a 2-column text file (Q [\AA{}\textasciicircum{}-1], S), linearly interpolated (0 outside).

Functions are given as strings evaluated at run time (MCViNE uses fparser): + - * / \textasciicircum{} (or **), \% , comparisons, \&\& ||, if(c,a,b), sin cos tan asin acos atan sinh cosh tanh exp log log10 log2 sqrt abs floor ceil int sign cbrt, pow atan2 min max hypot fmod, constants pi and e.

\textbf{Transport} (same as MCViNE HomogeneousNeutronScatterer): on the first passage the neutron is forced to scatter at a uniformly chosen depth, weighted by path length, attenuation exp(-(mu+sigma)x) and the scattering coefficient; on exit it is attenuated by exp(-(mu+sigma)L\_out). With order\textgreater{}1 further scatterings are sampled analogically (truncated exponential). With p\_transmit\textgreater{}0 that fraction of events is kept as the attenuated direct beam. mu(v) = absorption\_coefficient*2200/v. Events for which the kernel cannot scatter (kinematically forbidden) are absorbed.

Geometry: box (xwidth,yheight,zdepth), cylinder (radius,yheight), hollow cylinder (+thickness) or sphere (radius only). All vectors (Q, targets, reciprocal vectors, atom positions) are in the component's local frame.

Uses the shared runtime share/mcvine-lib.h/.c.

Example: MCViNE\_SQ(SQ="1+sin(Q*3)/(Q*3)", Qmin=0, Qmax=20, scattering\_coefficient=10, radius=0.01, yheight=0.05)

\subsection*{Input parameters}
Parameters in \textbf{boldface} are required; the others are optional.

\begin{longtable}{p{0.22\textwidth}p{0.12\textwidth}p{0.46\textwidth}p{0.14\textwidth}}
\toprule
\textbf{Name} & \textbf{Unit} & \textbf{Description} & \textbf{Default} \\
\midrule
\endhead
SQ & str & S(Q) expression (ignored if SQ\_file given) & "1" \\
SQ\_file & str & 2-column file Q S & 0 \\
Qmin & \AA{}$^{-1}$ & Lower Q bound & 0 \\
Qmax & \AA{}$^{-1}$ & Upper Q bound & 100 \\
absorption\_coefficient & m$^{-1}$ & Absorption coefficient at 2200 m/s (MCViNE convention; scales as 1/v) & 0 \\
scattering\_coefficient & m$^{-1}$ & Scattering coefficient & 0 \\
sigma\_abs & barn & Alternative: absorption cross section per unit cell at 2200 m/s (used when Vc\textgreater{}0) & 0 \\
sigma\_scat & barn & Alternative: scattering cross section per unit cell (used when Vc\textgreater{}0) & 0 \\
Vc & \AA{}$^{3}$ & Unit cell volume; when \textgreater{}0, coefficients are computed from sigma\_abs, sigma\_scat & 0 \\
radius & m & Outer radius of a cylinder (with yheight) or of a sphere (yheight=0) & 0 \\
xwidth & m & Width of a box sample & 0 \\
yheight & m & Height of a box or cylinder sample & 0 \\
zdepth & m & Depth of a box sample & 0 \\
thickness & m & Wall thickness of a hollow cylinder (0: filled) & 0 \\
pack & 1 & Packing factor (scales absorption and scattering coefficients) & 1 \\
p\_transmit & 1 & Monte Carlo fraction of events kept as unscattered transmitted beam (0: always scatter) & 0 \\
order & 1 & Maximum number of scattering events per neutron (1: single scattering) & 1 \\
\bottomrule
\end{longtable}

\subsection*{Links}
\begin{itemize}
  \item Component source code found in file \texttt{MCViNE\_SQ.comp}.
  \item MCViNE documentation: https://mcvine.github.io
\end{itemize}
\IfFileExists{contrib/MCViNE_SQ_static.tex}{\input{contrib/MCViNE_SQ_static.tex}}{}